% ===================================================================
%  সারণি নং 7 — শীতে গ্রাম। — বসন্তে খামারবাড়ি।
%
%  Ceci n'est PAS une transcription : c'est une traduction, faite en
%  2026, en bengali d'aujourd'hui. Voir texto/bn/10-sarani-01.tex
%  pour la consigne, qui vaut pour les seize fichiers.
% ===================================================================

%%K t07-noto-1 noto
\VUnotes{143.2mm}{%
(*) খাবারের বিশদ বিবরণের জন্য সারণি 6 ও 13 দেখুন।%
}

%%K t07-apar-1 apar
\VUsaut{4.0mm}
\VUcentre{13.2pt}{90}{দ্বিতীয় পর্যায়}
\VUsaut{3.0mm}
\VUfilet{16mm}
\VUsaut{5.6mm}
\VUtitre{15.0pt}{60}{সারণি নং 7}
\VUsaut{3.0mm}

%%K t07-tit-1 sub
\VUcentre{12.0pt}{0}{\textit{প্রথম দৃশ্য।}}
\VUsaut{3.0mm}

%%K t07-tit-2 sub
\VUpk{10.2pt}{40}{শীতে গ্রাম।}
\VUsaut{4.0mm}

%%K t07-c1-01-1 p
1. --- নতুন বছরের ছুটিতে আমি বাবার সঙ্গে নয়াগ্রামে গিয়েছিলাম, একটি ছোট গ্রাম
যেখানে তাঁর কিছু কাজ সারার ছিল।

%%K t07-c1-02-1 p
2. --- আমরা সেই \VUgras{ডাকগাড়ি} \textsuperscript{(11)} নিয়েছিলাম যা শহর ও
সেই গ্রামের মধ্যে চলে; কিন্তু ঠান্ডা খুব ছিল, আর এত অস্বস্তিকর গাড়িতে
যাত্রাটি একটুও সুখকর হয়নি।

%%K t07-c1-03-1 p
3. --- তাই আমাদের সত্যিই আনন্দ হল যখন দেখলাম \VUgras{গাড়োয়ান} চত্বরে
\textit{শ্বেত ভালুক} \VUgras{সরাইয়ের} \textsuperscript{(2)} সামনে গাড়ি
থামাচ্ছে। \VUgras{সরাইওয়ালা} \textsuperscript{(4)} নিজের দোরগোড়ায় আমাদের
অপেক্ষায় ছিল, শান্তভাবে নিজের \VUgras{কলকে} টানতে টানতে।

%%K t07-c1-03-2 p
সে আমাদের ভিতরে ডাকল, আর চুল্লিতে চটচট করে জ্বলা লতার আগুনের সামনে বসতে আমাকে
দুবার বলতে হল না। তারপর সেই কাঁপুনি-ধরানো \VUgras{উত্তুরে হাওয়ার} আমার আর কোনো
চিন্তা রইল না, যা পুরনো \VUgras{সাইনবোর্ড} \textsuperscript{(3)} কড়কড় করাত ও
চিমনি থেকে বেরোনো \VUgras{ধোঁয়া} \textsuperscript{(1)} কালো ঘূর্ণিতে উড়িয়ে
নিয়ে যেত।

%%K t07-c1-04-1 p
4. --- দুপুরের খাবারের সময় হয়েছিল। আর দেরি না করে আমরা সেই সাদাসিধা অথচ
সুস্বাদু খাবারের পুরো সদ্ব্যবহার করলাম যা আমাদের পরিবেশন করা হল (*)।

%%K t07-tit-3 sub
\VUpk{10.2pt}{40}{কয়েকটি সাক্ষাৎ।}
\VUsaut{3.0mm}

%%K t07-c1-05-1 p
5. --- খাওয়ার পরে আমি বাবার সঙ্গে গেলাম, যিনি বিমার প্রতিনিধি, তাঁর
খদ্দেরদের সঙ্গে দেখা করতে। প্রথমে আমরা \VUgras{কামারের} \textsuperscript{(5)}
কাছে গেলাম, যে সরাইয়ের পাশেই থাকে। সে একটি ঘোড়ার নাল ঠুকতে লেগে ছিল। আমি
তার সহকারীর জোর দেখে অবাক হলাম, যে ঘোড়ার \VUgras{খুর} নিজের চামড়ার
এপ্রনে শক্ত করে ধরে ছিল, যতক্ষণ কামার ভারী হাতুড়ির ঘায়ে \VUgras{নাল}
\textsuperscript{(6)} বসাতে থাকল।

%%K t07-c1-06-1 p
6. --- তারপর আমরা গ্রামের \VUgras{মুচি} \textsuperscript{(7)} বুড়ো জ্যাকবের
কাছে গেলাম। যদিও তার সাইনবোর্ডে একটি চমৎকার \VUgras{বুট} আঁকা, সে পুরনো
ক্ষয়ে যাওয়া জুতোয় নতুন সুকতলা লাগিয়েই কাজ চালাচ্ছিল।

%%K t07-c1-06-2 p
বুড়ো হেসে বলল: « শহরে আমি ছিলাম “জুতো-প্রস্তুতকারক” আর আমার খদ্দেরই ছিল না।
এখানে আমি “মুচি” আর কাজের কোনো অভাব নেই। »

%%K t07-c1-07-1 p
7. --- সে আমাদের নিজের প্রতিবেশী \VUgras{নাপিতের} \textsuperscript{(8)} কথা
বলল, যার খদ্দেররা সবাই অসন্তুষ্ট থাকে। তার \VUgras{ক্ষুর} সব সময় ভোঁতা থাকে
আর তার \VUgras{কাঁচি} কখনও ঠিক কাটে না। কিন্তু যেহেতু সে-ই গ্রামের একমাত্র
নাপিত, লোকের কাছে তার হাতেই দাড়ি কামানো ও চুল কাটানো ছাড়া উপায় নেই। সে
তার উপরে কৃপণ ও রুক্ষও।

%%K t07-c1-08-1 p
8. --- সৌভাগ্যক্রমে অন্য \VUgras{প্রতিবেশীরা} তার মতো নয়। \VUgras{গাড়ি
নির্মাতা} \textsuperscript{(9)}, উদাহরণস্বরূপ, ভালো মানুষ, পরিশ্রমী ও উপকারী।
তার কাছে গাড়ি মেরামতে দেওয়া সুখের কথা। তার কাজ সব সময় পরিপাটি হয়।

%%K t07-c1-09-1 p
9. --- বুড়ো জ্যাকবের \VUgras{জিনকারের} \textsuperscript{(10)} বিরুদ্ধেও কোনো
অভিযোগ ছিল না, যাকে সে কাজের চাপে কখনও কখনও \VUgras{সাজসরঞ্জাম} সেলাইয়ে
সাহায্য করে।

%%K t07-c1-09-2 p
আমার বাবা তাকে তার পরিবারের কথা জিজ্ঞেস করলেন। « সবাই ভালো আছে। আমার নাতনি
\VUgras{পাঠশালায়} \textsuperscript{(12)}। তার \VUgras{শিক্ষয়িত্রী}
\textsuperscript{(13)} তার উপরে খুব খুশি; সে ভালো \VUgras{ছাত্রী}
\textsuperscript{(14)}। সে আমার সেই দুষ্টু ছেলের মতো নয়; তার খেলা ছাড়া কিছু
মাথায় আসে না। \VUgras{প্রধান শিক্ষক} \textsuperscript{(15)} তাকে কিছুই শেখাতে
পারেন না। »

%%K t07-tit-4 sub
\VUsaut{3.0mm}
\VUpk{10.2pt}{40}{গ্রামের কথা।}
\VUsaut{3.0mm}

%%K t07-c1-10-1 p
10. --- তারপর বুড়ো আমাদের গ্রামের খবর শোনাল। একটি নতুন পাঠশালা হওয়ার কথা,
কারণ \VUgras{গ্রামভবনেরটি} খুব ছোট হয়ে পড়েছিল। এটি সহজ হত, কারণ
\VUgras{চক্কিওয়ালার} বাগানের একটি অংশ কিনে নেওয়াই যথেষ্ট ছিল। বেচারার
পক্ষে তা খুব ভালোই হত। ঠান্ডা পড়ার পর থেকে \VUgras{চক্কির}
\textsuperscript{(16)} \VUgras{পাট} গম পিষতে পারেনি, কারণ \VUgras{চাকা}
\textsuperscript{(17)} \VUgras{বরফে} \textsuperscript{(25)} আটকে গিয়েছিল, যা
\VUgras{নালার} \textsuperscript{(23)}।

%%K t07-c1-11-1 p
11. --- « নির্মাণের কথা যখন উঠল --- আমাদের নতুন \VUgras{গির্জা}
\textsuperscript{(18)} দেখেছেন? বরফের চাদরের নিচে সেটি বড় মনোরম লাগে।
\VUgras{কাকেরা} \textsuperscript{(19)}, যারা নিজেদের পুরনো ঘণ্টাঘর আর চিনতেই
পারে না, বিষণ্ণ হয়ে \VUgras{কবরস্থানের} \textsuperscript{(20)} বড় গাছটিতে বসে
থাকে। »

%%K t07-c1-12-1 p
12. --- কবরস্থানের নাম উঠতেই মুচির সেই লোকদের কথা মনে পড়ল যাদের আমার বাবা
চিনতেন এবং যারা গত বছর চলে গেছেন। « কত \VUgras{কবর} \textsuperscript{(21)}, আমার
ভালো মহাশয়, \VUgras{সাইপ্রেসের} \textsuperscript{(22)} নিচে আরও যোগ হল! মৃত্যু
আমাদের বুড়ো মুহুরিকে নিয়ে গেল। সারা গ্রাম তাঁর \VUgras{অন্ত্যেষ্টিতে}
ছিল। »

%%K t07-c1-13-1 p
13. --- « ওই যে তাঁর উত্তরসূরি, নালার উপরে বরফে পিছলে বেড়াচ্ছেন। একটু আগেই
এসেছিলেন যাতে আমি তাঁর \VUgras{বরফ-জুতোর} \textsuperscript{(26)} ফিতে জুড়ে দিই,
যা ছিঁড়ে গিয়েছিল। »

%%K t07-c1-14-1 p
14. --- « আর ওই যে, নালার ধারে, নতুন \VUgras{গ্রামরক্ষী}
\textsuperscript{(27)}। সে পুরনো সিপাহি আর গ্রামের দুষ্টু ছেলেদের যম। তার
\VUgras{পটি} \textsuperscript{(29)} আর তার \VUgras{ঝালরওয়ালা টুপি}
\textsuperscript{(28)} দেখলেই তারা পালিয়ে যায়। »

%%K t07-c1-15-1 p
15. --- « আরে! ওই যে বুড়ো জোসেফ, রুটিওয়ালার জন্য \VUgras{কাঠের আঁটির
গাড়ি} \textsuperscript{(30)} আনছে। --- ঠিক বলেছ, আমার বাবা বললেন, আমি তার
\VUgras{খচ্চরের} \textsuperscript{(31)} জোড়া চিনি। তার সঙ্গে আমার একটি কথা
আছে, এই তো সুযোগ। নমস্কার, বুড়ো জ্যাকব। »

%%K t07-c1-16-1 p
16. --- আমরা বেরোতেই যাচ্ছিলাম, এমন সময় গ্রামের ছেলেমেয়েরা পাঠশালা থেকে ছুটি
পেয়ে \VUgras{বরফের পিছল-পথে} \textsuperscript{(33)} ঝাঁপিয়ে পড়ল, এই ঝুঁকি
নিয়ে যে \VUgras{পড়ে গেলে} \textsuperscript{(34)} হাত-পা ভাঙবে। তাদের একজন
গাড়িওয়ালার ওখান থেকে নিজের \VUgras{স্লেজ} \textsuperscript{(32)} আনল, তাতে
নিজের এক সঙ্গীকে বসাল, আর দৌড়ে চলল।

%%K t07-c1-17-1 p
17. --- অন্যেরা একটি \VUgras{বরফের স্তূপের} \textsuperscript{(37)} দিকে ছুটল, যা
\VUgras{সেতুর} \textsuperscript{(40)} পাশে ছিল, আর সেই \VUgras{বরফের মানুষটিকে}
\textsuperscript{(39)} ঢিল ছুড়তে লাগল যা তারা কাল বানিয়েছিল, আর যাকে তারা
একটি টুপি, একটি কলকে ও কাঁধে ঝোলানো একটি ব্যাগ পরিয়েছিল।

%%K t07-c1-18-1 p
18. --- তাদের বরফশীতল হাওয়া ও ঠান্ডার পরোয়া ছিল না। তাদের অধিকাংশের
\VUgras{গলাবন্ধ} \textsuperscript{(35)} পর্যন্ত ছিল না, আর \VUgras{বরফের
গোলার} \textsuperscript{(38)} লড়াইয়ের পরে গরম হতে মুঠোভরা তপ্ত
\VUgras{চেস্টনাটই} \textsuperscript{(42)} তাদের যথেষ্ট ছিল, যা তারা বুড়ি
জ্যাকলিন \VUgras{বিক্রেতার} কাছে কিনত, যে নিজের খুব পাতলা \VUgras{চাদরে}
\textsuperscript{(43)} ঠান্ডায় কাঁপতে থাকে।

%%K t07-tit-5 sub
\VUsaut{3.0mm}
\VUcentre{12.0pt}{0}{\textit{দ্বিতীয় দৃশ্য।}}
\VUsaut{3.0mm}

%%K t07-tit-6 sub
\VUpk{10.2pt}{40}{বসন্তে খামারবাড়ি।}
\VUsaut{4.0mm}

%%K t07-c2-01-1 p
1. --- শীতের সব সুখ সত্ত্বেও \VUgras{বসন্তের} ফিরে আসাকে আমরা গভীর আনন্দে
স্বাগত জানাই।

%%K t07-c2-01-2 p
মে মাসে সন্ধ্যায় খামারের উঠান কী উৎসাহভরা ছবি গড়ে!

%%K t07-c2-02-1 p
2. --- দিগন্তে \VUgras{সূর্য} \textsuperscript{(1)} ধীরে ধীরে ঢলছে, আর
\VUgras{মাঠকে} \textsuperscript{(3)} নিজের লাল \VUgras{রশ্মিতে} \textsuperscript{(2)}
ভরিয়ে দিচ্ছে। পরিশ্রমী \VUgras{হালচাষি} \textsuperscript{(6)} নিজের
\VUgras{খেত} \textsuperscript{(4)} চষার তাড়ায়।

%%K t07-c2-02-2 p
নিজের \VUgras{লাঙল} \textsuperscript{(5)} দিয়ে, যা একটি শক্তিশালী
\VUgras{ঘোড়া} \textsuperscript{(7)} টানে, সে \VUgras{সিতা} \textsuperscript{(8)}
টানে, যাতে \VUgras{বপনকারী} \textsuperscript{(9)} মুঠো মুঠো সেই \VUgras{বীজ}
ছড়াবে যা থেকে সোনালি শিষ বেরোবে।

%%K t07-c2-03-1 p
3. --- নিজেদের কাস্তে দিয়ে \VUgras{ঘাসকাটিয়েরা} \textsuperscript{(11)} মাঠের
ঘাস কেটে ফেলেছে, শুকানেওয়ালারা \VUgras{শুকনো ঘাস} \textsuperscript{(10)}
\VUgras{আঁচড়া} \textsuperscript{(12)} ও খুরপি দিয়ে উলটেপালটে শুকিয়ে নিয়েছে,
আর এখন তারা ফিরছে।

%%K t07-c2-03-2 p
কেউ কেউ বলদে টানা ভারী গাড়ির আগে আগে চলছে, কাঁধে নিজেদের যন্ত্রপাতি নিয়ে।
অন্যেরা, বেশি অলস, সুগন্ধি শুকনো ঘাসে আরাম করে গা এলিয়ে দিয়েছে।

%%K t07-c2-04-1 p
4. --- চলতে চলতে তারা \VUgras{মালিকে} \textsuperscript{(16)} বন্ধুত্বের সঙ্গে
অভিবাদন জানায়, যে \VUgras{বাগানে} \textsuperscript{(13)} লেগে আছে,
\VUgras{ফলগাছের} ডালছাঁটাই করছে ও \VUgras{নাশপাতিগাছে} \textsuperscript{(14)}
কলম বসাচ্ছে। \VUgras{আলুবোখারার গাছ} \textsuperscript{(15)} ফুলে ভরা, আর ভালো
সার পাওয়া \VUgras{সবজিবাগানে} \textsuperscript{(17)} তরকারি, বাঁধাকপি ও
লেটুস আগে থেকেই ভালো হচ্ছে।

%%K t07-c2-04-2 p
নিচু দেয়ালের উপরে একটি সুন্দর \VUgras{ময়ূর} \textsuperscript{(18)} নিজের
চমৎকার পালক দেখাতে পেখম মেলে।

%%K t07-tit-7 sub
\VUpk{10.2pt}{40}{খামারের উঠান।}
\VUsaut{3.0mm}

%%K t07-c2-05-1 p
5. --- \VUgras{আস্তাবলের} \textsuperscript{(19)} দরজা খোলা, আর জাবপাত্র ও
\VUgras{বিছানার খড়} \textsuperscript{(23)} দেখা যাচ্ছে সেই \VUgras{ঘোটকীর}
\textsuperscript{(21)}, যাকে \VUgras{সহিস} \textsuperscript{(20)} এইমাত্র খুলেছে।
\VUgras{বাচ্চা ঘোড়া} \textsuperscript{(22)}, নিজের লম্বা পায়ে এত মজার, লাফাতে
লাফাতে মায়ের পিছনে চলে।

%%K t07-c2-06-1 p
6. --- \VUgras{কুয়োর} \textsuperscript{(29)} পাশে \VUgras{গরুরা}
\textsuperscript{(25)} সেই কাঠের \VUgras{গামলা} \textsuperscript{(26)} থেকে জল খেতে
যায় যা তাদের ঘাট। \VUgras{বাছুর} \textsuperscript{(24)} এই ফাঁকেই দুধ খেয়ে
নেয়, আর \VUgras{বলদ} \textsuperscript{(27)}, খড়ের ভালো গন্ধ পেয়ে, আনন্দে
হাম্বা ডাকে ও নিজের বলবান \VUgras{শিংওয়ালা} \textsuperscript{(28)} মাথা গর্বে
তোলে।

%%K t07-c2-07-1 p
7. --- সবাই কাজে লেগে। \VUgras{কাজের মেয়ে} \textsuperscript{(32)} কুয়োর
\VUgras{দড়ি} \textsuperscript{(31)} ধরে আছে। সে গরুদের খাওয়ানোর জন্য একটি
\VUgras{বালতিতে} \textsuperscript{(30)} জল তুলছে। \VUgras{খামারঘরের}
\textsuperscript{(33)} পাশে একজন লোক ছড়ানো খড় জড়ো করছে, আর \VUgras{খামারবাড়ির}
\textsuperscript{(34)} দরজার সামনে এক বুড়ি \VUgras{সুতাকাটুনি}
\textsuperscript{(35)}, নিজের চরকা ও নিজের \VUgras{টাকু} নিয়ে, পুরনো দিনের মতো
\VUgras{শণ} বা \VUgras{পাট} কাটছে।

%%K t07-tit-8 sub
\VUsaut{3.0mm}
\VUpk{10.2pt}{40}{খামারের মালিক।}
\VUsaut{3.0mm}

%%K t07-c2-08-1 p
8. --- \VUgras{খামারের মালিক} \textsuperscript{(36)} এই সব কাজ চালান ও নিজের
লোকদের হুকুম দেন। তিনি তাদের বলেন \VUgras{মই} \textsuperscript{(40)} ও
\VUgras{কোদাল} \textsuperscript{(39)} ঢেকে রাখতে, যা সেই \VUgras{ঘরের}
\textsuperscript{(38)} পাশে ফেলে রাখা হয়েছে যা \VUgras{কুকুরের}
\textsuperscript{(37)}।

%%K t07-c2-09-1 p
9. --- তাঁর হাঁকে একটি \VUgras{পায়রা} \textsuperscript{(41)} চমকে ওঠে, যা পুরো
গতিতে \VUgras{পায়রাখোপের} \textsuperscript{(42)} দিকে উড়ে যায়, আর
\VUgras{মুরগিরাও} \textsuperscript{(43)}, যারা তাড়াতাড়ি \VUgras{খোপে}
\textsuperscript{(44)} ফিরে যায়।

%%K t07-c2-10-1 p
10. --- \VUgras{ভেড়াশালার} \textsuperscript{(48)} সামনে, এই যে বুড়ো
\VUgras{রাখাল} \textsuperscript{(45)} নিজের \VUgras{পাল} \textsuperscript{(49)}
ফিরিয়ে আনছে। নিজের \VUgras{লাঠি} \textsuperscript{(46)} ধরে, যা তার সঙ্গ ততটাই
ছাড়ে না যতটা তার ফ্যাকাশে \VUgras{কুর্তা} \textsuperscript{(47)}, সে খোঁয়াড়ের
দিকে \VUgras{মেষ} \textsuperscript{(50)}, \VUgras{ভেড়া} \textsuperscript{(53)},
\VUgras{মেষশাবক} \textsuperscript{(54)}, \VUgras{ছাগল} \textsuperscript{(51)} ও
\VUgras{ছাগলছানা} \textsuperscript{(52)} হাঁকিয়ে নেয়। তার বিশ্বস্ত
\VUgras{কুকুর} \textsuperscript{(55)} আর্গুস সবার পিছনে আসে ও ঘেউ ঘেউ করে
পিছিয়ে-পড়াদের এগিয়ে দেয়।

%%K t07-c2-11-1 p
11. --- এই সব হইচই সেই ভালো \VUgras{দুধেল গাইয়ের} \textsuperscript{(56)} শান্তি
নষ্ট করে না, যে নিজের \VUgras{ঘণ্টি} \textsuperscript{(57)} বাজাতে থাকে, যখন তাকে
\VUgras{গোয়ালঘরের} \textsuperscript{(58)} দরজার সামনে দোয়ানো হয়।

%%K t07-c2-12-1 p
12. --- সে যেন বোঝে যে তাকে নিজের দুধ দিতেই হবে, যাতে \VUgras{খামারের গিন্নি}
\textsuperscript{(60)} \VUgras{মন্থনদণ্ডে} \textsuperscript{(61)} \VUgras{মাখন}
বানাতে পারেন, বা সেই চমৎকার \VUgras{পনির} \textsuperscript{(64)} যা
\VUgras{গামলার} \textsuperscript{(63)} উপরে রাখা তক্তায় নিংড়ে যাচ্ছে।

%%K t07-c2-13-1 p
13. --- সারা \VUgras{দুগ্ধশালা} \textsuperscript{(59)} \VUgras{খামারের মজুরের}
\textsuperscript{(65)} কল্যাণে খুব পরিষ্কার থাকে, যে এইমাত্র সেই বড় বালতিতে
\VUgras{দুধের কৌটো} \textsuperscript{(62)} ধুয়েছে যা সে তুলে ধরে আছে।

%%K t07-tit-9 sub
\VUsaut{3.0mm}
\VUpk{10.2pt}{40}{মুরগিশালা।}
\VUsaut{3.0mm}

%%K t07-c2-14-1 p
14. --- নিজের ভারী খড়মে সে কিছুটা বেঢপ ভঙ্গিতে হাঁটে, আর তাতে
\VUgras{টার্কি} \textsuperscript{(68)}, \VUgras{হাঁসিরা} \textsuperscript{(66)},
\VUgras{হাঁসেরা} \textsuperscript{(67)} ও \VUgras{রাজহাঁস} \textsuperscript{(69)}
উড়ে যায়, যারা নিজেদের \VUgras{ঠোঁট} \textsuperscript{(70)} চওড়া করে খুলে
অস্থির কোলাহল করে।

%%K t07-c2-14-2 p
\VUgras{মোরগ} \textsuperscript{(71)}, যে বেশি লড়াকু, নিজের লাল \VUgras{ঝুঁটি}
\textsuperscript{(72)} তোলে, নিজের \VUgras{লেজের} \textsuperscript{(73)} পালক
ঝাড়ে আর একটি গমগমে « কোঁকর-কোঁ » ছাড়ে।

%%K t07-c2-15-1 p
15. --- একটি \VUgras{মা-মুরগি} \textsuperscript{(74)} \VUgras{গোবরের স্তূপে}
\textsuperscript{(81)} খুঁটে খাচ্ছে, নিজের \VUgras{ছানাদের} \textsuperscript{(75)}
সঙ্গে। \VUgras{শুয়োরছানা} \textsuperscript{(78)} নিজের মায়ের দিকে ছুটে যায়,
সেই বিশাল \VUgras{শুয়োরির} \textsuperscript{(77)} দিকে যাকে আমরা খোঁয়াড়ের
পাশে দেখতে পাই।

%%K t07-c2-16-1 p
16. --- নিজেদের খাঁচায় \VUgras{খরগোশেরা} \textsuperscript{(79)} সেই কচি
\VUgras{ঘাস} \textsuperscript{(80)} সাগ্রহে খাচ্ছে যা মালিকের মেয়ে তাদের এনে
দেয়। জেনি নিজের খরগোশদের খুব ভালোবাসে, আর প্রতিদিন, খোপ থেকে টাটকা
\VUgras{ডিম} \textsuperscript{(76)} আনার পরে, সে নিজের ছোট্ট বন্ধুদের সঙ্গে
দেখা করতে আসে।
\VUsaut{6.0mm}
\VUfilet{22mm}
